% META: {"id": "1NSI-WEB-ME-003", "chapitre": "1NSI-WEB-IHM", "type_objet": "methode", "capacites": ["P-IHM-03A", "P-IHM-03B", "P-IHM-03C"], "capacites_codes": ["C4", "C5", "C6"], "methodes": ["M3"], "mode_creation": "ex_nihilo", "sources_inspiration": [], "status": "needs_review"}
% BEGIN-VERIFY
% cookies_client = {}
% def serveur_definit_cookie(nom, valeur):
%     cookies_client[nom] = valeur
% def client_envoie_requete_avec_cookies(url):
%     return {"url": url, "cookies": dict(cookies_client)}
% avant = client_envoie_requete_avec_cookies("/accueil")
% assert avant["cookies"] == {}
% serveur_definit_cookie("panier", "2 articles")
% apres = client_envoie_requete_avec_cookies("/paiement")
% assert apres["cookies"] == {"panier": "2 articles"}
% assert apres["url"] == "/paiement"
% assert avant["cookies"] == {}
% END-VERIFY
\begin{fichemethode}{M3}{Situer ce qui s'exécute côté client, ce qui est mémorisé, ce qui est chiffré}

\textbf{Quand l'utiliser ?} Dès qu'un énoncé décrit un échange entre un navigateur et un
serveur : ordre des étapes, rôle des cookies, ou question sur la confidentialité.

\textbf{Pas à pas :}
\begin{enumerate}
  \item \textbf{Trace deux colonnes}, client et serveur, et place chaque action dans la sienne.
        Le \textbf{client} est le navigateur de l'utilisateur ; le \textbf{serveur} est la
        machine distante. Rien ne s'exécute simultanément des deux côtés : une requête part,
        une réponse revient.
  \item \textbf{Ordonne les étapes.} Une action côté serveur ne peut pas précéder la requête
        qui la déclenche. Si l'énoncé demande « dans quel ordre », c'est cette contrainte qui
        donne la réponse.
  \item \textbf{Demande-toi ce qui survit à la réponse.} Le serveur ne se souvient de rien
        entre deux requêtes. Ce qui est mémorisé l'est \textbf{dans le client}, sous forme de
        cookie, et \textbf{renvoyé automatiquement} à chaque requête suivante vers ce site.
  \item \textbf{Pour la confidentialité, sépare trois moments} : avant l'envoi (dans le
        navigateur), pendant la transmission (sur le réseau), après réception (sur le serveur).
        HTTPS chiffre la \textbf{transmission} — et seulement elle.
\end{enumerate}

\textbf{Exemple rédigé.} Un visiteur charge \lstinline{/accueil}, ajoute un article, puis
charge \lstinline{/paiement}.

\begin{python}
avant = client_envoie_requete_avec_cookies("/accueil")
serveur_definit_cookie("panier", "2 articles")
apres = client_envoie_requete_avec_cookies("/paiement")
\end{python}

La requête vers \lstinline{/accueil} ne transporte \textbf{aucun} cookie : il n'en existe pas
encore. Après que le serveur a défini le cookie, la requête vers \lstinline{/paiement}
transporte \lstinline{{"panier": "2 articles"}} — sans que le visiteur ait à le redemander.
Le serveur n'a rien mémorisé : c'est le client qui a rapporté l'information.

\textbf{Erreur fréquente.} Croire que HTTPS rend les données invisibles partout. Le serveur
destinataire, lui, lit les données en clair : il faut bien qu'il les traite. HTTPS empêche un
tiers de les lire \emph{sur le réseau}, pas le destinataire de les connaître.

\end{fichemethode}
